\chapter*{About this record}
\addcontentsline{toc}{chapter}{About this record}

This is the complete record of the \cytofair{} project from its first registered commit on 22 August 2026 to
the end of the follow-up experiments on 20 October 2026. It is deliberately exhaustive. It keeps the original
Study~1 record intact as Part~II, and adds everything that happened afterwards: the review that turned Study~1's
weaknesses into a scaling roadmap, the separately preregistered Study~2, the engineering needed to run \RecAllTasks
GPU tasks on a shared host with a four-core CPU allowance, every registered result, the eight follow-up
hypotheses of Amendment~8, the bugs and setbacks, and what the project established, refuted, and left open.

\begin{keybox}{Evidence boundary}
Every number in Parts~I and III--VIII is generated from committed outputs: the registered analysis files in
\artifact{results/study2/analysis} and \artifact{results/study2/followup}, the exported run records in
\artifact{results/study2/run_records}, the power simulation, the provenance maps, and the git history. The
generators are \artifact{scripts/study2/paper_numbers.py}, \artifact{record/scripts/make_record_tables.py}, and
\artifact{record/scripts/make_record_figures.py}; no result is typed by hand. Part~II is the Study~1 record of 13
September 2026, carried over verbatim (labels prefixed so they cannot collide), and its numbers come from the Study~1
result files as they stood then.
\end{keybox}

\begin{warningbox}{What ``preregistered'' means here}
Study~2 and Amendment~8 were registered as commits in version control before any model whose result they govern was
trained, and the registration commit's hash is recorded inside every run record. That is strong internal evidence, but
the commits have not yet been pushed or deposited with a third party, so an outside reader cannot yet verify their
dates independently. Amendment~8 was written after Study~2's results were known and is reported as a registered
follow-up, not as part of Study~2.
\end{warningbox}

\section*{How the record is organised}

\begin{description}[leftmargin=2.2em,style=nextline]
\item[Part I, the lifecycle at a glance] One chapter with the whole story: timeline, every registered test and its
verdict, the hypothesis map, and the answers to ``what did we get out of it?''.
\item[Part II, Study 1] The frozen Study~1 record: motivation, provenance recovery, structural leakage, literal overlap,
the parent-identity probe, the SIPaKMeD and Cx22 rescoring, statistics, execution failures, verification, and release.
\item[Part III, from Study 1 to Study 2] The roadmap and our response to each item, Phase~0 checks, the Study~2
preregistration, power analysis, every design (splits, factorial, transplants, adversary, CRIC), the structural closed
forms, and the post hoc reanalysis of Study~1.
\item[Part IV, engineering and execution] The host, its constraints, the GPU-resident input path, determinism, the
trainer's artifact schema, the scheduler, failure handling, smoke tests, the execution timeline, compute, archiving,
and the complete log of errors and fixes.
\item[Part V, Study 2 results] H1 to H5 with every sensitivity analysis and secondary analysis, training curves,
calibration, confusion matrices, and per-architecture detail.
\item[Part VI, Amendment 8] The follow-up question, the eight registered hypotheses, their designs, the preparer bug
(Amendment~9), every result, and the mechanism by elimination.
\item[Part VII, synthesis] What was established, refuted, and left open; what worked and what did not; lessons;
threats to validity; recommendations; publication options; decisions only the authors can make.
\item[Part VIII, appendices] The Study~1 appendices, complete run ledgers for every Study~2 and follow-up run, variance
components, power grid, risk report, the full commit chronology, registration hashes, and a reproduction cookbook.
\end{description}

\section*{Reading paths}

\begin{itemize}[leftmargin=1.6em]
\item \emph{Ten minutes:} Part~I.
\item \emph{The science:} Part~I, Chapter~\ref{r:ch:roadmap}, Part~V, Part~VI, Chapter~\ref{r:ch:synthesis}.
\item \emph{Reproducing or extending the work:} Part~IV, Appendix~\ref{r:app:cookbook}, and the ledgers.
\item \emph{Writing a paper from it:} Chapter~\ref{r:ch:synthesis} (claims and their strength) and
Section~\ref{r:sec:publication}.
\end{itemize}

\section*{Conventions}

P1 is a stratified crop-level 80:20 split; P1-prime (P1$'$) a size-matched control; P2 a parent-disjoint split (field,
source image, or slide); P3 a parent-disjoint split after merging near-duplicate-linked parents. ``pp'' means percentage
points. Unless stated otherwise, intervals are 95\%, times are Asia/Kolkata (IST), and ``supported'' means a registered
test passed its registered criterion (adjusted $p<0.05$ with the predicted sign). In figures, crop-level (leaky) series
are orange, parent-level blue, and component-level aqua; leak models are orange and clean models blue.
